\section*{Bab \ref{ch:b2:mammal-reproduction} --- Perkembangbiakan Seksual Mamalia}

\begin{solution}{exo:b2:mammal-reproduction:1}
Spermatogonium membelah sejak akil balig seumur hidup; oogonium hanya
pada janinnya. Satu \omterm{def:b2:meiosis-heredity:meiosis}{meiosis} memberi empat sperma tetapi hanya satu sel
telur (dan tiga badan kutub). Sebuah sperma memakan kira-kira 64 hari
ditambah dua pekan pematangan; sebutir sel telur terhenti pada \omterm{prop:b2:meiosis-heredity:stages}{profase
I} sejak sebelum lahir sampai bulan \omterm{def:b2:mammal-reproduction:ovary}{ovulasinya} --- berpuluh tahun ---
lalu terhenti lagi pada metafase II sampai pembuahannya. Seorang pria
membuat $10^{8}$ sperma sehari, sekitar $10^{12}$ seumur hidup; seorang
perempuan mengovulasikan kira-kira 400 sel telur.
\end{solution}

\begin{solution}{exo:b2:mammal-reproduction:2}
Pengikatan pada \omterm{prop:b2:mammal-reproduction:fertilisation}{zona pelusidanya}; \omterm{prop:b2:mammal-reproduction:fertilisation}{reaksi akrosom} dan pencernaan sebuah
jalan; peleburan membrannya dan masuknya nukleus spermanya; gelombang
kalsium dan \omterm{prop:b2:mammal-reproduction:fertilisation}{reaksi korteks} (pengerasan zonanya, hilangnya reseptornya:
penghadang polispermi); penyelesaian \omterm{def:b2:meiosis-heredity:meiosis}{meiosis} II dan pengeluaran badan
kutub kedua; pembentukan kedua pronukleusnya, replikasi DNA-nya,
mitosis pertamanya.
\end{solution}

\begin{solution}{exo:b2:mammal-reproduction:3}
\omterm{def:b2:mammal-reproduction:implantation}{Trofoblas} (lapisan luarnya) membuat \omterm{def:b2:mammal-reproduction:placenta}{plasenta} dan selaputnya; \omterm{def:b2:mammal-reproduction:implantation}{massa sel
dalamnya} membuat lembaganya. \omterm{def:b2:mammal-reproduction:implantation}{Implantasinya} mulai pada hari ke-6 sampai
ke-7 sesudah pembuahan, di dalam pelapis rahim yang telah disiapkan
progesteron dari \omterm{def:b2:mammal-reproduction:ovary}{korpus luteumnya}.
\end{solution}

\begin{solution}{exo:b2:mammal-reproduction:4}
Oksigen: difusi. Glukosa: pengangkutan terbantu. Sel darah merah
induknya: tidak menyeberang (kedua darahnya tidak pernah bercampur).
IgG: pengangkutan berperantara reseptor. Urea: difusi, dari janinnya ke
induknya. Alkohol: berdifusi bebas. Bakteri: sebagian besar tidak
menyeberang (sebagian, seperti sifilis dan listeria, menyeberang).
\end{solution}

\begin{solution}{exo:b2:mammal-reproduction:5}
$m = 3, 1, 0.4, 0.1$: $P = 0.95, 0.63, 0.33, 0.095$. Di bawah separuh
ketika $m < \ln 2 = 0.69$, yakni di bawah $7\times 10^{7}$ sperma.
Kurvanya menjenuh: di atas kira-kira $10^{8}$, sperma tambahan hampir
tidak menambah apa pun, dan di bawahnya kesuburannya turun nyaris
sebanding dengan cacahnya.
\end{solution}

\begin{solution}{exo:b2:mammal-reproduction:6}
Sebelum akil balig: $7\times 10^{5}$ hilang dalam 13 tahun (\num{4750}
hari): kira-kira 150 tiap hari. Sesudahnya: \num{299000} hilang dalam
37 tahun (\num{13500} hari): 22 tiap hari; \omterm{def:b2:mammal-reproduction:ovary}{ovulasinya}
$400/\num{13500} = 0.03$ tiap hari. Atresianya melampaui \omterm{def:b2:mammal-reproduction:ovary}{ovulasinya}
tujuh ratus banding satu; persediaannya habis bukan karena diovulasikan
melainkan karena mati.
\end{solution}

\begin{solution}{exo:b2:mammal-reproduction:7}
Dewasanya: $(30/27)^{2.7} = 1.33$, $S = 1.33/2.33 = 0.57$. Janinnya:
$(30/19)^{2.7} = 3.4$, $S = 3.4/4.4 = 0.77$. Pada tekanan rendah
\omterm{def:b2:mammal-reproduction:placenta}{plasentanya}, darah janinnya memuat oksigen seperlima lebih banyak
daripada darah orang dewasa, lalu melepaskannya di dalam jaringan pada
tekanan yang membuat hemoglobin dewasa sudah separuh kosong.
\end{solution}

\begin{solution}{exo:b2:mammal-reproduction:8}
$2^{k} = 25$: $k = 4.6$ kali pelipatduaan, yaitu \qty{9.3}{d}: hari
ke-16 sampai ke-17 sesudah pembuahan; hari ke-30 sampai ke-31 sesudah
haid terakhirnya --- kira-kira hari haidnya terlewat.
\end{solution}

\begin{solution}{exo:b2:mammal-reproduction:9}
$3\times 10^{-8}\times 6\times 10^{4}\times 20/10^{-3} =
\qty{36}{mL/min}$ terhadap permintaan \qty{25}{mL/min}: marjinnya
turun dari delapan menjadi satu setengah; janinnya menjadi kekurangan
oksigen di bawah beban tambahan apa pun, tumbuh lambat, dan mungkin
perlu dilahirkan lebih awal.
\end{solution}

\begin{solution}{exo:b2:mammal-reproduction:10}
(a) \omterm{def:b2:mammal-reproduction:testis}{Testisnya} terbentuk lalu menyekresikan kedua hormonnya;
AMH-nya menghapus saluran Müllernya, tetapi testosteronnya tak dapat bekerja:
saluran Wolffnya menyusut dan alat kelamin luarnya betina --- seorang
perempuan dengan testis dan tanpa rahim. (b) Testosteronnya bekerja,
AMH-nya tidak: seorang jantan dengan epididimis, vas deferens, dan alat
kelamin jantan, ditambah sebuah rahim dan saluran telur. (c) Ovarium
dan turunan Müllernya (tanpa testis, tanpa AMH); androgen adrenalnya
menjantankan alat kelamin luarnya sampai taraf yang berbeda-beda. (d)
\emph{SRY} membuat testis, yang memaksakan jalan jantannya: seorang
pria dengan dua kromosom X, mandul karena gen Y \omterm{def:b2:mammal-reproduction:testis}{spermatogenesisnya}
tidak ada.
\end{solution}

\begin{solution}{exo:b2:mammal-reproduction:11}
Yang paling lazim (\qty{68}{\%}) adalah pembelahan hari ke-4 sampai
ke-8: sesudah \omterm{def:b2:mammal-reproduction:implantation}{trofoblasnya} terbentuk (sehingga satu \omterm{def:b2:mammal-reproduction:placenta}{plasenta}) tetapi
sebelum amnionnya (sehingga dua kantung). Sebelum hari ke-3 tiap
belahannya membuat \omterm{def:b2:mammal-reproduction:implantation}{trofoblasnya} sendiri, sehingga dua \omterm{def:b2:mammal-reproduction:placenta}{plasenta}
(\qty{30}{\%}); sesudah hari ke-8 amnionnya sudah diletakkan dan
dibagi bersama (\qty{2}{\%}). Membelah \omterm{def:b2:mammal-reproduction:implantation}{massa sel dalamnya} sesudah
\omterm{def:b2:mammal-reproduction:implantation}{trofoblasnya} terpancang jelas merupakan peristiwa yang paling mudah.
\end{solution}

\begin{solution}{exo:b2:mammal-reproduction:12}
Seekor marsupial menanamkan sedikit sebelum kelahirannya: kehamilan
yang pendek, seekor anak yang mungil, dan sebuah laktasi yang dapat
dipotong pendek dengan menjatuhkan anak kantungnya kalau kekeringan
datang, sehingga risiko induknya tiap upayanya rendah dan ia dapat
mencoba lagi seketika. Seekor mamalia berplasenta menanamkan banyak
pada kehamilan yang panjang dengan seekor anak yang besar dan
berkembang baik yang bertahan hidup lebih baik saat lahir, dengan
ongkos seekor induk yang terikat berbulan-bulan dan rentan selama itu.
Di lingkungan yang produktif dan teramalkan, kebertahanan anak mamalia
berplasenta yang lebih tinggi menang; iklim Australia yang tak menentu
mengganjar pilihan marsupial untuk meninggalkan anaknya, dan
keterpencilannya menahan mamalia berplasenta di luar sampai baru-baru
ini.
\end{solution}

\begin{solution}{pb:b2:mammal-reproduction:1}
\textbf{1.} $m = 3\times 10^{8}\times 10^{-8} = 3$; $P = 1 - e^{-3} =
0.95$.
\textbf{2.} $1 - e^{-3}(1 + 3) = 0.80$: tanpa penghadangnya, empat dari
lima pembuahannya akan polispermik dan lembaganya triploid.
\textbf{3.} $0.15/5\times 10^{-5} = \qty{3000}{s}$, yaitu lima puluh
menit; kerutan rahim dan saluran telurnya yang membawanya.
\textbf{4.} $300/3\times 10^{8} = 10^{-6}$: hilang di dalam vagina yang
asam, di dalam lendir serviksnya, kepada fagosit di dalam rahimnya, dan
di dalam saluran telur yang keliru.
\textbf{5.} Dari tiga hari sebelum \omterm{def:b2:mammal-reproduction:ovary}{ovulasi} sampai satu hari
sesudahnya: kira-kira empat hari.
\textbf{6.} $m = 0.5$, $P = 0.39$; harapan cacah daurnya $1/P = 2.5$.
\textbf{7.} Mitokondria spermanya yang sedikit, di dalam bagian
tengahnya, ditandai lalu dimusnahkan sesudah masuk; sel telurnya
mengusung seratus ribu.
\textbf{8.} $120/16 = 7.5$: tujuh pembelahan, $2^{7} = 128$ sel ---
kira-kira seratus.
\textbf{9.} $\log_2 25 = 4.6$ kali pelipatduaan, yaitu \qty{9.3}{d}:
hari ke-16 sampai ke-17; hari ke-30 sampai ke-31 sesudah haidnya.
\textbf{10.} \omterm{def:b2:mammal-reproduction:ovary}{Korpus luteumnya} menyusut kira-kira dua belas hari sesudah
\omterm{def:b2:mammal-reproduction:ovary}{ovulasi} kecuali kalau hCG menyelamatkannya; tanpanya progesteronnya
turun, pelapisnya luruh, dan lembaganya hilang bersama haidnya.
\textbf{11.} $\log_2 10^{5} = 16.6$ kali pelipatduaan; pada pekan ke-10
\omterm{def:b2:mammal-reproduction:placenta}{plasentanya} membuat progesteronnya sendiri dan \omterm{def:b2:mammal-reproduction:ovary}{korpus luteumnya} tidak
lagi dibutuhkan.
\textbf{12.} Sel telur duduk pada \omterm{prop:b2:meiosis-heredity:stages}{profase I} berpuluh tahun dan kohesin
yang menahan bivalennya melapuk, sehingga kegagalan pemisahan pada
\omterm{def:b2:meiosis-heredity:meiosis}{meiosis} I naik tajam dengan umur induknya; sperma dibuat segar dalam
dua bulan.
\textbf{13.} Sesudah \omterm{def:b2:mammal-reproduction:implantation}{trofoblasnya} (hari ke-5) dan sebelum amnionnya
(hari ke-8): satu \omterm{def:b2:mammal-reproduction:placenta}{plasenta}, dua kantung.
\textbf{14.} $3\times 10^{-8}\times 1.2\times 10^{5}\times 20/3.5\times
10^{-4} = \qty{206}{mL/min}$.
\textbf{15.} $7\times 3.5 = \qty{24.5}{mL/min}$: marjin delapan kali.
\textbf{16.} $500\times 0.16 = \qty{80}{mL/min}$; penyariannya
$24.5/80 = \qty{31}{\%}$.
\textbf{17.} $5\times 3.5\times 1440 = \qty{25}{g/d}$; termuat di dalam
\qty{28}{L} darah induknya; aliran hariannya \qty{720}{L}: janinnya
mengambil kira-kira \qty{4}{\%} glukosa yang lewat.
\textbf{18.} Hemoglobin janinnya \qty{77}{\%} jenuh pada
\qty{30}{mmHg}, janinnya memiliki lebih banyak sel darah merah dan
keluaran jantung yang tinggi, dan jaringannya teradaptasi bekerja pada
tekanan rendah --- ia hidup, pada dasarnya, di puncak Everest.
\textbf{19.} Pertukaran gas tanpa paru-paru yang bekerja, nutrisi dan
pengeluaran tanpa usus atau ginjal yang bekerja, serta hormon yang
mempertahankan kehamilannya; ia tidak dapat menahan tiap racun ---
alkohol, nikotin, dan sebagian virus menyeberang.
\textbf{20.} Regangan serviks $\to$ hipotalamus $\to$ oksitosin $\to$
kerutan $\to$ regangan yang lebih besar, sampai kelahirannya membuang
rangsangannya. Sebelum cukup bulan, rahim yang didominasi progesteron
memiliki sedikit reseptor oksitosin dan tanpa sambungan renggang,
sehingga kerutannya tidak menyebar dan gelungnya tidak dapat menutup.
\textbf{21.} $750\times 2.9 = \qty{2.2}{MJ}$; ongkosnya $2.2/0.8 =
\qty{2.7}{MJ}$.
\textbf{22.} Pertumbuhannya $30\times 6 = \qty{0.18}{MJ}$,
pemeliharaannya $300\times 4 = \qty{1.2}{MJ}$: seluruhnya
\qty{1.4}{MJ}, kira-kira dua pertiga energi susunya; sisanya masuk ke
kegiatan, kehangatan, dan kehilangan.
\textbf{23.} $270\times 2.7 = \qty{730}{MJ}$, yaitu dua setengah kali
kehamilannya.
\textbf{24.} Mengisap menaikkan prolaktin lalu menekan pelepasan
berdenyut hormon yang menggerakkan \omterm{def:b2:mammal-reproduction:ovary}{ovulasi}, sehingga \omterm{def:b2:mammal-reproduction:ovary}{ovulasinya} baru
kembali berbulan-bulan kemudian; tanpa kontrasepsi hal ini
merenggangkan kelahirannya dua sampai tiga tahun.
\textbf{25.} $P = 0.95$; ujinya positif pada hari ke-17 sesudah
pembuahan (hari ke-31 daurnya); marjin oksigennya delapan kali; susu
sehari memakan ongkos \qty{2.7}{MJ} bagi induknya.
\end{solution}
